\documentclass[10pt,a4paper]{amsart}
\usepackage[margin=19mm]{geometry}
\usepackage{amsmath,amssymb,mathtools}
\usepackage[scr=rsfs,cal=euler]{mathalfa}
\usepackage{xfrac}
\usepackage[unicode,hidelinks,bookmarks=false]{hyperref}
\newcommand{\ie}{i.e.\ }
\newcommand{\cf}{cf.\ }
\newcommand{\resp}{respectively\ }
\newcommand{\Subsection}{\subsection}
\providecommand{\nobreakdash}{\nobreak\hbox{-}}
\setlength{\emergencystretch}{2em}
\multlinegap0pt
\newtheorem{prop}{Proposition}
\newcommand{\cR}{\mathcal{R}}
\title{On monomial resonance}
\author{C\u{a}lin Spiridon}
\date{Source: arXiv:2606.19006v1 (CC BY 4.0)}
\begin{document}
\maketitle

\noindent\emph{Source and licence.} On monomial resonance. Version arXiv:2606.19006v1 (CC BY 4.0), distributed by arXiv under the Creative Commons Attribution 4.0 International licence, http://creativecommons.org/licenses/by/4.0/, as recorded in the arXiv licence metadata for this exact version and shown on the abstract page https://arxiv.org/abs/2606.19006v1. Full licence notice: \texttt{licenses/arxiv-2606.19006-CC-BY-4.0.txt}.\par\smallskip

\noindent\textit{Editorial scope.} Counted unit: the article's prop prop\_terminal\_vertex (source0.tex lines 507--509). The article's complete original proof of it is delivered with it (source0.tex lines 511--540, one \texttt{proof} environment of 30 lines). No mathematics is written, repaired or re-derived: the statement and the proof are the article's own lines, in the article's own notation, and no retained line is substituted or re-flowed.  Retained uncounted as the source-local context the proof needs: lines 476--478.  Nothing was omitted from any retained span, so the package carries no omission record.  No retained line contains a citation, so the wrapper carries no bibliography and no bibliography entry.  No retained line contains a figure environment, an image-inclusion command or drawing code, so no image is delivered and the package declares no asset dependency.  Editorial formatting only, in addition: an AMS \texttt{amsart} preamble that restates the article's own declarations and macros used by the retained lines, namely the environments \texttt{prop} on separate counters, together with the standard packages those lines need.  The retained environments print in this document as Proposition 1 in order of appearance, whereas the article prints the same items with the numbering of its own full document.  The complete native source of the article as distributed in the arXiv e-print for this exact version is bundled unchanged as \texttt{sources/203107/source0.tex}.
\begin{prop} \label{prop_triangle}
If the graph $\Gamma$ contains a triangle (i.e. a cycle on $3$ vertices), there exists a non-monomial subspace $K$ such that $\cR_\Gamma = \cR(V,K)$.
\end{prop}

\begin{prop} \label{prop_terminal_vertex}
If the graph $\Gamma$ has a terminal vertex, there exists a non-monomial subspace $K$ such that $\cR_\Gamma = \cR(V,K)$.
\end{prop}

\begin{proof}
Without loss of generality, suppose that $1$ is a terminal vertex and $(1,2) \in \mathsf{E}$. If $\Gamma$ contains a triangle, the conclusion already follows from Proposition \ref{prop_triangle}, so we may also assume that, for instance, $(3,4) \notin \mathsf{E}$.

Define $K$ to be the subspace of $\bigwedge^2V$ generated by $e_1 \wedge e_2 + e_3 \wedge e_4$ and all the totally decomposable vectors $e_i \wedge e_j$ such that $(i,j) \in \mathsf{E} \setminus \{(1,2)\}$. Equivalently, $K^\perp$ is generated by $f_1 \wedge f_2 - f_3 \wedge f_4$ and all the totally decomposable vectors $f_i \wedge f_j$ such that $(i,j) \notin \mathsf{E}$.

Let $a  = \sum_{i=1}^n a_i f_i \in V^\vee$ be a non-zero vector. Note that $a \in \cR_\Gamma$ if and only if there exists $b = \sum_{i=1}^n b_i f_i$, which is non-proportional to $a$, such that
\begin{equation} \label{eq:graph_terminal}
    a_i b_j = a_j b_i, \ \forall \ (i,j) \in \mathsf{E}
\end{equation}
and $a \in \cR(V,K)$ if and only if there exists a vector $b' = \sum_{i=1}^n b'_i f_i$, which is non-proportional to $a$, such that
\begin{equation} \label{eq:non-monomial_terminal}
    \begin{cases}
        a_1 b'_2 - a_2 b'_1 = a_3 b'_4 - a_4 b'_3 \\
        a_i b'_j = a_j b'_i, \ \forall \ (i,j) \in \mathsf{E}\setminus \{(1,2)\}
    \end{cases}
\end{equation}
If $a_2 = 0$, then $a \in \cR_\Gamma \cap \cR(V,K)$, since we may take $b = b' = f_1$. Therefore, we assume that $a_2 \neq 0$.

Suppose that $a \in \cR_\Gamma$. Then, there exists a vector $b \in V^\vee$, such that $a \wedge b \neq 0$ and the components of $b$ with respect to the basis $\{f_1, \ldots, f_n\}$ satisfy the relations (\ref{eq:graph_terminal}). Let us take $b' = \sum_{i=1}^n b'_i f_i$, where $b'_i = b_i$, for all $i \ge 2$ and
\[
b'_1 = \frac{1}{a_2}(a_1 b_2 - a_3 b_4 + a_4 b_3)
\]
It is then straightforward to check that $a \wedge b' \neq 0$ and the components of $b'$ satisfy the relations (\ref{eq:non-monomial_terminal}). Therefore, $a \in \cR(V,K)$.

Suppose now that $a \in \cR(V,K)$. Then, there exists a vector $b' \in V^\vee$, such that $a \wedge b' \neq 0$ and the components of $b'$ with respect to the basis $\{f_1, \ldots, f_n\}$ satisfy the relations (\ref{eq:non-monomial_terminal}). Consider now $b = \sum_{i=1}^n b_i f_i$, where $b_i = b'_i$, for all $i \ge 2$ and
\[
b_1 = \frac{a_1 b_2}{a_2}
\]
Note that $a \wedge b \neq 0$ and the components of $b$ satisfy the relations (\ref{eq:graph_terminal}), hence $a \in \cR_\Gamma$.
\end{proof}

\end{document}
